% !TEX root = main.tex
% Included within Appendix D; all symbols below are local to the clean binary model.
\subsection{Clean distractors: learned suppression}
\label{app:J-clean}



This subsection proves the clean binary routing part of \cref{thm:joint}. Each prediction has the same $m$ designated distractors
with zero key and value; the only shared trainable routing parameter is their penalty $B$.
The binary start count is uniform and the match bit has probability $p$.
The supplied useful-site score model is
\begin{equation}
 \begin{gathered}
 \mathrm{score}(i_t)=0,\qquad \mathrm{score}(j_t)=(-1)^{b_t}a,\qquad
 \mathrm{score}(\text{distractor})=-B,\\
 r_b=\frac{1}{1+e^{(-1)^ba}+me^{-B}}.
 \end{gathered}
 \label{eq:J-clean-scores}
\end{equation}

Let $m\ge1$, $0<p<1/2$, $d_{\mathrm h}\ge1$, $\delta\in[0,1]$, and define
\[
 \begin{gathered}
 \mathcal D=me^{-B},\quad x=-a,\quad
 r_0=\frac1{1+e^{-x}+\mathcal D},\quad r_1=\frac1{1+e^x+\mathcal D},\\
 L=(1-p)\log(1+e^{-c-gr_0})+p\log(1+e^{c+gr_1}).
 \end{gathered}
\]
Let $E_0=(1-p)\sigma(-c-gr_0)$, $E_1=p\sigma(c+gr_1)$.
Cyclic symmetry gives $\mathbf U=\upsilon_0\mathbf I+\upsilon_1\mathbf R$ and $\mathbf W=w(\mathbf I-\mathbf R)$, with
$c=\initscale(\upsilon_0-\upsilon_1)$ and $g=2\initscale vw$. Thus the two correct-class margins are $c+gr_0$ and $-c-gr_1$.
The same binary readout metric as in part (i), now with the soft weights in \cref{eq:J-clean-scores}, gives the exact joint flow:
\begin{align*}
 \ddt{c}&=\initscale^2(E_0-E_1),&
 \ddt{g}&=\initscale^2v^2(E_0r_0-E_1r_1),\\
 a&=\frac{\initscale^2qk}{\sqrt{d_{\mathrm h}}},&
 \ddt{a}&=-\frac{\initscale^4(q^2+k^2)}{d_{\mathrm h}}\pd{L}{a},\\
 \pd{L}{a}&=g(E_0e^a r_0^2+E_1e^{-a}r_1^2),&
 \ddt{B}&=\nu g\mathcal D(E_0r_0^2-E_1r_1^2).
\end{align*}
The two useful-site biases are held at zero; all designated distractor biases equal $-B$.
Here $0<\nu\le1$. A Euclidean scalar penalty uses $\nu=1$; tying $m$ distractor bias entries to $-B$ under their Frobenius metric gives $\nu=1/m$.
Initialize $c=g=0$, $B=\delta$, $qk<0$, with
$\initscale/2\le |q(0)|,|k(0)|,|v|\le2\initscale$.
Put
\begin{equation}
 h^\star=\log\frac{1-p}{p},\quad \bar h^\star=\min\{h^\star,1\},\quad
 z=\frac{p\log2}{2},\quad
 b_\star=\frac1{2+e^{-1}}-\frac1{2+e},
\end{equation}
\begin{equation}
 T_F=\frac{16\log2}{b_\star^2p^2(1-e^{-z})^2},\qquad
 G_{\max}=2+4T_F,\qquad x_\star=\max\{1,\log(8(m+2)),\log(2G_{\max})\}.
 \label{eq:CR-constants}
\end{equation}
All these quantities are independent of $\initscale$; $e^{x_\star}=\max\{e,8(m+2),2G_{\max}\}$ is polynomially bounded in $m,p^{-1}$.

\paragraph{Explicit bound for \cref{thm:joint}(ii).}
Assume
\begin{equation}
 \begin{aligned}
 0<\initscale\le\initscale^\star_{\mathrm{clean}}:=\min\Biggl\{&\frac12,\frac1{2(m+2)},
 \left(\frac{\nu}{16(\log m+2)}\right)^{1/4},\\
 &\frac1{8e^{x_\star}\sqrt{d_{\mathrm h}}(m+2)(\log(4x_\star\sqrt{d_{\mathrm h}})+4)}\Biggr\}.
 \end{aligned}
 \label{eq:CR-scale}
\end{equation}
The simultaneous flow enters the region $x\ge x_\star$, $\mathcal D\le1$ in finite time. This region persists long enough for $L\le z$, after which both transitions have strictly positive margins forever. A sufficient total time is
\begin{align}
 T\le{}&\frac{h^\star}{2\initscale^2 p(1-p)(1-e^{-h^\star/4})}\notag\\
 &+\frac{48\sqrt{d_{\mathrm h}}(m+2)(1+e^{x_\star}+m)^2}{\initscale^4p\bar h^\star}
       \log\frac{4x_\star\sqrt{d_{\mathrm h}}}{\initscale^4}
 +\frac{32(m+2)}{\initscale^4p}+\frac{T_F}{\initscale^4}.\label{eq:CR-time}
\end{align}
Thus $T=\initscale^{-4}\operatorname{poly}(m,d_{\mathrm h},p^{-1},(\bar h^\star)^{-1},\log(1/\initscale))$.

\subsubsection{An invariant for readout growth and routing}
\begin{lemma}[Clean alignment and loss bounds]
\label{lem:CR-align}
For all time, $c\le h^\star$, $g\ge0$, and $x\ge x(0)>0$. Moreover $L\le\log2$. Whenever $g\le1$,
 $c\ge-2\log2-1,\qquad E_0\ge p/4,\qquad E_1\ge p/12.$ 
\end{lemma}
\begin{proof}
At $g=0$, $c\le h^\star$ implies $E_0\ge E_1$, and $r_0>r_1$ implies $\ddt{g}>0$.
At $c=h^\star$, $g>0$ gives $E_0<E_1$, hence $\ddt{c}<0$. Also $\pd{L}{a}>0$ when $g>0$, so $\ddt{x}>0$.
These boundary signs preserve the region from the initial state. Gradient flow decreases the loss.
For $g\le1$, $(1-p)\log(1+e^{-c-gr_0})\le\log2$ implies $c+gr_0\ge-2\log2$, hence the lower bound on $c$.
Since $c\le h^\star$,
$E_0\ge(1-p)/(1+e^{h^\star+1})\ge p/(e+1)>p/4$.
The lower bound on $c$ yields $E_1\ge p/(1+4e)>p/12$.
In coordinates $u$ with the constant metric whitened,
\[
 \int_0^T\|\ddt{u}\|^2\,\mathrm d t=L(0)-L(T)\le\log2,\qquad
 \|u(T)-u(0)\|\le\sqrt{T\log2}.
\]
The vector field is smooth on all finite parameter values. The continuation criterion of \citet[Corollary~2.15]{teschl2012ode} therefore gives global existence.
\end{proof}

Define
 $J=\log\frac{E_0r_0}{E_1r_1}.$ 
As long as $J\ge0$, both $g$ and $B$ are nondecreasing. Indeed
 $E_0r_0^2-E_1r_1^2 =r_0(E_0r_0-E_1r_1)+E_1r_1(r_0-r_1)\ge0.$ 
Furthermore,
 $\ddt{(B+2e^B/m)}=(1+2/\mathcal D)\ddt{B}\ge\nu\frac{g\ddt{g}}{\initscale^2v^2},$ 
because $(\mathcal D+2)r_0\ge1$. Thus, on any interval starting at zero on which $J\ge0$,
\begin{equation}
 B(t)+\frac{2e^{B(t)}}m\ge\delta+\frac{2e^\delta}m+\frac{\nu g(t)^2}{2\initscale^2v^2}.
 \label{eq:CR-potential}
\end{equation}
With $\mathcal D_0=me^{-\delta}\le m$, this implies
\begin{equation}
 \nu g^2\mathcal D^2\le\initscale^2v^2(m+2)^2.
 \label{eq:CR-budget}
\end{equation}
To see this, multiply \eqref{eq:CR-potential} by $\mathcal D^2$ and use
$\mathcal D^2\log(\mathcal D_0/\mathcal D)\le \mathcal D_0^2/(2e)$ and $\mathcal D\le \mathcal D_0$:
 $\nu g^2\mathcal D^2\le2\initscale^2v^2\{\mathcal D^2\log(\mathcal D_0/\mathcal D)+2\mathcal D\} \le\initscale^2v^2(\mathcal D_0^2/e+4\mathcal D_0)\le\initscale^2v^2(m+2)^2.$ 

\begin{lemma}[Monotone clean readout and routing]
\label{lem:CR-monotone}
On any interval starting at initialization on which either $g\le1$ or both $x\ge x_\star$ and $g\le G_{\max}$ at every time,
$J>0$ persists. Consequently $g$ and $B$ are nondecreasing throughout that interval.
\end{lemma}
\begin{proof}
Initially $J=h^\star+\log(r_0/r_1)>0$. Suppose $J$ has a first zero.
Write $\Delta r=r_0-r_1>0$, $\sigma_0=\sigma(c+gr_0)$, $\sigma_1=\sigma(-c-gr_1)$,
$\kappa_0=e^{-x}r_0$, $\kappa_1=e^xr_1$.
At $J=0$, $E_0r_0=E_1r_1$, so
 $\ddt{g}=0,\quad \ddt{c}=-\initscale^2 E_1\Delta r/r_0,\quad \ddt{B}=\nu g\mathcal DE_1r_1\Delta r.$ 
Differentiation gives
\begin{align*}
 \ddt{J}={}&-(\sigma_0+\sigma_1)\ddt{c}
 +[\kappa_0+\kappa_1-g(\sigma_0\kappa_0r_0-\sigma_1\kappa_1r_1)]\ddt{x}\\
 &+\mathcal D[\Delta r-g(\sigma_0r_0^2+\sigma_1r_1^2)]\ddt{B}.
\end{align*}
The count term dominates the possibly negative routing term: by \eqref{eq:CR-budget}, $r_0r_1\le1/4$, and $v^2(m+2)^2\le1$,
\begin{align*}
 &\initscale^2 E_1\Delta r(\sigma_0+\sigma_1)/r_0
 -\nu g^2\mathcal D^2E_1r_1\Delta r(\sigma_0r_0^2+\sigma_1r_1^2)\\
 &\hspace{1em}\ge\frac{E_1\Delta r(\sigma_0+\sigma_1)}{r_0}
       [\initscale^2-\nu g^2\mathcal D^2r_0r_1]\\
 &\hspace{1em}\ge\frac{3\initscale^2 E_1\Delta r(\sigma_0+\sigma_1)}{4r_0}>0.
\end{align*}
The remaining routing term $\mathcal D\Delta r\ddt{B}$ is nonnegative.
The coefficient of $\ddt{x}$ is nonnegative if $g\le1$.
For $x\ge x_\star$, $\kappa_1/\kappa_0=e^{2x}r_1/r_0\ge e^x\ge2G_{\max}$, so that coefficient is also positive for $g\le G_{\max}$.
Up to a first zero, \eqref{eq:CR-potential}--\eqref{eq:CR-budget} apply and give $\ddt{J}>0$, contradicting the first-zero condition $\ddt{J}\le0$.
\end{proof}

\begin{proof}[Proof of \cref{thm:joint}(ii)]
\emph{Basin entry.}
First, $x$ reaches $x_\star$ before $g$ reaches $1$. Until either event, $\mathcal D\le m$, $g$ is nondecreasing, and
\begin{align*}
 \frac1x\ddt{x}&\ge\frac{2\initscale^2g}{\sqrt{d_{\mathrm h}}}E_0e^{-x}r_0^2\\
 &\ge\frac{2ge^{-x_\star}}{\sqrt{d_{\mathrm h}}v^2(m+2)}\ddt{g}.
\end{align*} 
The first inequality uses $q^2+k^2\ge2\sqrt{d_{\mathrm h}}x/\initscale^2$.
Integration gives
 $\log\frac{x}{x(0)}\ge\frac{e^{-x_\star}g^2}{\sqrt{d_{\mathrm h}}v^2(m+2)}.$ 
If $g=1$ occurred first, this would contradict
 $v^2e^{x_\star}\sqrt{d_{\mathrm h}}(m+2)\log\frac{x_\star}{x(0)}\le\frac12.$ 
This inequality follows from the stated scale threshold, $x(0)\ge\initscale^4/(4\sqrt{d_{\mathrm h}})$, and
$\initscale^2\log(4x_\star\sqrt{d_{\mathrm h}}/\initscale^4)\le\initscale[\log(4x_\star\sqrt{d_{\mathrm h}})+4]$.

For a quantitative time bound, let $g_*=v^2\bar h^\star/[2(m+2)]$ and
$\mu_{\mathrm{prior}}=p(1-p)(1-e^{-h^\star/4})$.
Before $g$ reaches $g_*$, it satisfies $g<h^\star/4$. Whenever $c\le h^\star/2$,
 $E_0-E_1\ge(1-p)\sigma(-3h^\star/4)-p\sigma(3h^\star/4)\ge\mu_{\mathrm{prior}}.$ 
Thus $c$ would reach $h^\star/2$ by $T_0=h^\star/(2\initscale^2\mu_{\mathrm{prior}})$ if $g_*$ had not already been reached.
While $c$ increases,
 $\ddt{g}=v^2r_0\ddt{c}+\initscale^2v^2(r_0-r_1)E_1\ge\frac{v^2}{m+2}\ddt{c}.$ 
Consequently reaching $c=h^\star/2$ forces $g\ge v^2h^\star/[2(m+2)]\ge g_*$.
By monotonicity, $g\ge g_*$ thereafter, up to the relevant stopping event.

Until $x=x_\star$, $g\le1$ and $E_1\ge p/12$, hence
\[
 \frac{\ddt{x}}{x}\ge\frac{2\initscale^2g_*}{\sqrt{d_{\mathrm h}}}
       \frac{p}{12(1+e^{x_\star}+m)^2}
 \ge\frac{\initscale^4p\bar h^\star}{48\sqrt{d_{\mathrm h}}(m+2)(1+e^{x_\star}+m)^2}.
\]
Integrating this inequality after $T_0$ gives the second term of \cref{eq:CR-time}; if $x=x_\star$ is reached before $T_0$,
that earlier entry only shortens the bound.

After $x=x_\star$, if $\mathcal D>1$, \eqref{eq:CR-potential} implies
 $g^2\le\frac{2\initscale^2v^2}{\nu}(\log m+2)\le\frac12.$ 
So $g$ cannot reach $1$ before $\mathcal D\le1$. Also $r_1/r_0\le1/8$, $E_0\ge p/4$, $E_1\le p$, and
 $\ddt{g}\ge\frac{\initscale^2v^2 p}{8(m+2)}.$ 
Thus $\mathcal D\le1$ is reached within another $8(m+2)/(\initscale^2v^2 p)\le32(m+2)/(\initscale^4p)$.
If $\mathcal D\le1$ was reached earlier, monotonicity retains it while $x$ grows to $x_\star$.

\emph{Completion and retention.}
In the basin, $r_0-r_1\ge b_\star$: the contrast increases with $x$ and decreases with $\mathcal D$, so its minimum for $x\ge1$, $\mathcal D\le1$
is the value defining $b_\star$. The readout gradient vector is
$\bigl(E_1-E_0,E_1r_1-E_0r_0\bigr)$.
Its two-by-two coefficient matrix has determinant magnitude $r_0-r_1$ and Frobenius norm at most $2$.
The smallest squared singular value is therefore at least $(r_0-r_1)^2/4$, and
$\min\{\initscale^2,\initscale^2v^2\}\ge\initscale^4/4$ gives
 $\ddt{L}\le-\frac{\initscale^4b_\star^2}{16}(E_0^2+E_1^2).$ 
While $L\ge z$, at least one example has loss at least $z$, hence
$E_0^2+E_1^2\ge p^2(1-e^{-z})^2$. Therefore $L\le z$ occurs within $T_F\initscale^{-4}$ time as long as the basin persists.
At entry $g<1$. Since $|\ddt{g}|\le\initscale^2v^2\le4\initscale^4$, throughout that completion interval
$g\le1+4T_F<G_{\max}$. \Cref{lem:CR-monotone} for $x\ge x_\star$, $g\le G_{\max}$ therefore preserves both $x\ge x_\star$ and $\mathcal D\le1$ throughout the interval.
Any nonpositive correct-class margin contributes at least $p\log2=2z$ to the loss. Thus $L\le z$ gives positive margins for both transitions, and loss monotonicity preserves them forever.
\end{proof}

\input{independent_routing_proof}
